\chapter{导论}\label{ch:1}
\srcpage{5}
本章概述从整体类域论通向几何 Langlands 理论的整体、局部和量子形式的路径。关键在于，Galois 数据与自守函数之间的经典关系具有几何提升：将函数替换为模叠上的层，将 Frobenius 特征值替换为 Hecke 特征同构。现代的函子化表述见文献\cite{GR24a,AG15}。

本讲义的安排如下。\chref{2}单独介绍 $\Bun_G$ 及其基本的叠论描述和阿代尔描述；同一章还表述 Hecke 特征性质，并给出 $\GL_n$ 的微小 Hecke 论证。\chref{3}讨论 Whittaker 模型，\chref{4}构造并分析 Laumon 层。第\ref{ch:5}–\ref{ch:9}章，尤其是\chref{8}中的构造，证明主要的 Fourier 变换、洁净延拓和消失结论。范畴形式在\chref{10}、\chref{11}、\chref{12}和\chref{13}中展开；Hitchin 理论和局部方面的输入见\chref{14}和\chref{15}；整体证明则安排在\chref{16}。

\section{几何类域论}\label{sec:1.1}
\index{类域论!几何类域论|hyperpage}\index{阿代尔|hyperpage}\index{非分歧|hyperpage}
设 $K$ 为整体域，$\A_K$ 为其阿代尔环，$\Gal_{\overline K/K}$ 为其绝对 Galois 群。整体类域论给出一个“几乎”同构；在本概述中，这个词有意使用得不甚精确，后续讲授将以严格陈述取代这一概括：
\[
\theta_K:K^\times\backslash\A_K^\times\longrightarrow
\Gal_{\overline K/K}^{\mathrm{ab}}\simeq\Gal_{K^{\mathrm{ab}}/K},
\]
其中 $K^{\mathrm{ab}}\subset\overline K$ 是最大阿贝尔扩张。这个同构由局部 Frobenius 元刻画：对每个有限素位 $v$，
\[\theta_K(\pi_v)=\Fr_v.\]
这里 $\pi_v$ 是 $K_v$ 的一个素元所确定的类，模 $\cO_v^\times$ 后有良好定义；$\Fr_v$ 是 $\Gal_{\overline{\kappa_v}/\kappa_v}$ 中几何 Frobenius 的一个提升，模惯性群 $I_v$ 后有良好定义。

\begin{transnote}{含糊点拨}
「几乎」在此是数学限定而非修辞。互反映射 $\theta_K$ 对数域只是单射、像稠密而不满；取 $\pi_0(K^\times\backslash\A_K^\times)$ 后核恰为单位分支，此时才得到同构（Artin 互反律的连通核即此分支）。对 $\mathbb F_q$ 上的函数域——本节随后转入的情形——精确陈述把 $\Gal^{\mathrm{ab}}$ 换成 Weil 群 $W_K^{\mathrm{ab}}$，$\theta_K$ 才是同构。下一段关于 $\mathbb Q_\ell^\times$ 完全不连通、因而无法检测连通核的说明，正是这一点的直接后果。
[来源：nLab "class field theory"，https://ncatlab.org/nlab/show/class+field+theory（Last revised 2022-10-18；2026-11-17 访问）；并见原书 p.5 第 2–3 段及 p.5 对 Weil 群的说明。]
\end{transnote}

\begin{theorem}{非分歧类域论：特征标形式}{1.1.1}
设 $\ell\ne\operatorname{char}(K)$。存在一个几乎双射
\[
\{\text{非分歧特征标 }\xi:K^\times\backslash\A_K^\times\to\Qell^\times\}
\longleftrightarrow
\{\text{非分歧特征标 }\rho:\Gal_{\overline K/K}\to\Qell^\times\},
\]
它由等式
\[\xi(\pi_v)=\rho(\Fr_v)\]
刻画。非分歧是指：对所有有限素位 $v$，$\xi$ 在 $\cO_v^\times$ 上平凡，$\rho$ 在 $I_v$ 上平凡。
\end{theorem}

对数域而言，由于 $\Qell^\times$ 完全不连通，互反映射的连通核无法被取值于 $\Qell^\times$ 的特征标检测到。对 $\mathbb F_q$ 上的函数域而言，精确表述使用 Weil 群 $W_{\overline K/K}$。

现在假设 $K$ 是 $k=\mathbb F_q$ 上一条光滑完备曲线 $X$ 的函数域。记 $\cO_K^\times:=\prod_v\cO_v^\times$。商 $K^\times\backslash\A_K^\times/\cO_K^\times$ 是离散的，Weil 定理给出
\[
K^\times\backslash\A_K^\times/\cO_K^\times\simeq\Pic(X),
\qquad\pi_v\longmapsto\cO_X(v).
\]
因此，非分歧自守特征标与 Picard 群的特征标相同：
\[
\{\text{非分歧 }\xi:K^\times\backslash\A_K^\times\to\Qell^\times\}
\simeq\{\xi:\Pic(X)\to\Qell^\times\}.
\]

\Needspace{4\baselineskip}
在 Galois 一侧，对于 $\Spec K$ 的一个几何点 $\eta$，有 $\Gal_{\overline K/K}\simeq\pi_1(\Spec K,\eta)$。非分歧表示经由以下满射分解：
\[\pi_1(\Spec K,\eta)\longrightarrow\pi_1(X,\eta)\]
\srcpage{6}
等价地，
\[
\begin{aligned}
\{\text{非分歧 }\rho:\Gal_{\overline K/K}\to\Qell^\times\}
&\simeq\{\rho:\pi_1(X,\eta)\to\Qell^\times\}\\
&\simeq\{X\text{ 上秩为一的 }\ell\text{-进局部系统 }\sigma\}.
\end{aligned}
\]
于是，类域论成为几乎双射
\[
\{\xi:\Pic(X)\to\Qell^\times\}
\longleftrightarrow\{X\text{ 上秩为一的 }\ell\text{-进局部系统 }\sigma\}.
\]

剩下的任务是将等式 $\xi(\pi_v)=\rho(\Fr_v)$ 用几何方式表达出来。设 $Y$ 是 $k=\mathbb F_q$ 上的代数簇，$\cF$ 是 $Y$ 上的 $\ell$-进层。对 $y\in Y(k)$，茎 $\cF_{\bar y}$ 带有 $\Gal_{\bar k/k}$ 的作用；定义其迹函数为
\[f_{\cF}(y):=\Tr(\Fr_y,\cF_{\bar y}).\]
这就是 Grothendieck 的层–函数对应。若 $\rho$ 对应局部系统 $\sigma$，则在 $k$-点 $v\in X(k)$ 处，
\[\rho(\Fr_v)=f_\sigma(v)=\Tr(\Fr_v,\sigma_{\bar v}).\]
对闭点也使用相同的记号，Frobenius 则取相应的几何 Frobenius。

\begin{corollary}{非分歧特征标关系的几何形式}{1.1.2}
存在几乎双射
\[
\{\xi:\Pic(X)\to\Qell^\times\}\longleftrightarrow
\{X\text{ 上秩为一的 }\ell\text{-进局部系统 }\sigma\},
\]
由 $\xi(\cO_X(v))=f_\sigma(v)$ 刻画。
\end{corollary}

现在将 Picard 群替换为相对 Picard 概形。对每个仿射测试概形 $S/k$，
\[
\Pic_{X/k}(S):=\Pic(X\times_k S)/\Pic(S),\qquad
\Pic_{X/k}(k)=\Pic(X).
\]
我们希望在 $\Pic_{X/k}$ 上得到一个几何对象 $\Aut_\sigma$，其迹函数就是与 $\sigma$ 相联系的特征标：
\origdiagram{p006-diagram.pdf}{0.8971}
对 $\cF$ 应施加的正确几何条件，是 Hecke 特征性质。

\begin{definition}{$\GL_1$ 的 Hecke 特征层}{1.1.3}
设 $\cF$ 是 $\Pic_{X/k}$ 上的层，$\sigma$ 是 $X$ 上秩为一的局部系统。若 $\cF$ 配备同构
\[\operatorname{add}^*\cF\simeq\sigma\boxtimes\cF,\]
且该同构满足连续两次修改的通常交换性条件，就称 $\cF$ 是以 $\sigma$ 为特征值的 Hecke 特征层。这里
\[
\operatorname{add}:X\times_k\Pic_{X/k}\longrightarrow\Pic_{X/k},
\qquad(v,L)\longmapsto L(v).
\]
\end{definition}

\begin{transnote}{含糊点拨}
「连续两次修改的通常交换性条件」指把 $\sigma$-等变性做成对有序点组一致的相容数据。对 $v\ne w$，先改 $v$ 再改 $w$ 与先改 $w$ 再改 $v$ 在 $\Pic_{X/k}$ 上落到同一点（$L(v)(w)=L(w)(v)$）；条件要求由此得到的两个同构 $\operatorname{add}^*_{v,w}\cF\simeq\sigma_v\otimes\sigma_w\otimes\cF$ 相等。缺了这条一致性，$\cF$ 只在每条 $\operatorname{add}$ 的纤维上「是特征的」，拼不成 $X\times_k\Pic_{X/k}$ 上带对称群相容性的层，\lemref{1.1.4} 也就无从谈起。一般 $G$ 的完整形式见 \S\ref{sec:2.6}：Hecke 函子给出的同构须与 $\Rep(\checkG)$ 中的张量积相容。
[依据：原书 p.6 定义 1.1.3 的措辞与 \S\ref{sec:2.6} 的一般定义；余圈条件本身出自 [BD91] 对 Hecke 特征层的原始构造。]
\end{transnote}

\Needspace{8\baselineskip}
\begin{lemma}{}{1.1.4}
设 $\cF$ 是以 $\sigma$ 为特征值的 Hecke 特征层，并假设 $\cF$ 在单位点 $\cO_X\in\Pic_{X/k}$ 处的限制是常值层，从而相应函数将单位元映为 $1$。在 $\Pic(X)=\Pic_{X/k}(k)$ 上令 $\xi=f_{\cF}$。则 $\xi$ 是特征标，且对每个素位 $v$，
\[f_{\cF}(\cO_X(v))=f_\sigma(v).\]
\end{lemma}
\srcpage{7}
\lemref{1.1.4}说明，Hecke 特征同构是在层的层面上对类域论中乘法性方程和 Frobenius 特征值方程的替代。

\begin{theorem}{Deligne}{1.1.5}
对 $X$ 上每个秩为一的局部系统 $\sigma$，在 $\Pic_{X/k}$ 上都存在本质唯一的 Hecke 特征层 $\Aut_\sigma$，其特征值为 $\sigma$；而且 $\Aut_\sigma$ 是局部系统。这里“本质唯一”是指，忽略与 $\Shv(\Spec k)$ 中对象的张量积后唯一；这种张量积保持特征值不变。
\end{theorem}

\thref{1.1.5}是纯几何陈述，Deligne 的证明也是几何的。因此，它就是非分歧整体几何类域论，等价地说，是 $\GL_1$ 的非分歧整体几何 Langlands 对应。

\section{非分歧整体几何 Langlands 对应}\label{sec:1.2}
\index{Langlands 对应!整体几何|hyperpage}\index{Hecke 特征层|hyperpage}\index{局部系统|hyperpage}
现在从 $\GL_1$ 转到 $\GL_n$，再转到 $k$ 上的分裂约化群 $G$。对 $X$ 上非常值约化群概形的版本也很重要，但不属于这组入门讲授的内容。

从现在起，$X$ 是完美域 $k$ 上光滑、完备且几何连通的曲线，$K=K(X)$ 是它的函数域。固定一种合理的层理论，其系数域 $e$ 代数闭。对性质良好的 $k$-概形或叠 $Y$，以 $\Shv_c(Y)^\heartsuit$ 表示阿贝尔范畴，以 $\Shv_c(Y)$ 表示导出范畴。标准选择为：
\begin{itemize}
\item Betti 情形：$k=\mathbb C$，使用复拓扑下的反常层；
\item $\ell$-进情形：$k$ 为有限域，使用反常 $\ell$-进 Weil 层；
\item de Rham 情形：$k$ 为特征 $0$ 的代数闭域，使用凝聚右 $D$ 模。
\end{itemize}

\begin{transnote}{易错预警}
上标 $\heartsuit$ 是 $t$-结构理论中的标准记号，指相应 $t$-结构的\textbf{心}（heart），不是「阿贝尔群值的层」。此处的「阿贝尔范畴」指所选层理论本身给出的对象范畴（Betti 情形即反常层范畴），而 $\Shv_c(Y)$ 是它的导出范畴。于是 $\cF\in\Shv_c(Y)^\heartsuit$ 与把 $\cF$ 当作导出范畴中的复形不可混同：书中的 Hecke 函子、Fourier 变换与上同调都在 $\Shv_c(Y)$ 层面计算，只在 $t$-正合性成立时才能把陈述降到心层面。\conjref{1.3.1} 与 \thref{1.3.2} 的差别正是这一处：前者是心范畴之间的 $t$-正合等价，后者放弃 $t$-正合性。
[来源：nLab "t-structure"，https://ncatlab.org/nlab/show/t-structure（Last revised 2023-04-20；2026-11-17 访问）；nLab "perverse sheaf"，https://ncatlab.org/nlab/show/perverse+sheaf（Last revised 2026-10-16；2026-11-17 访问）；并见原书 p.7 \S\ref{sec:1.2} 开头的约定。]
\end{transnote}

Galois 一侧的推广是形式上的：
\[
\begin{array}{c|l}
\GL_n&X\text{ 上秩为 }n\text{ 的局部系统 }\sigma,\\\hline
G&X\text{ 上的 }G\text{-局部系统 }\sigma.
\end{array}
\]
对于与 $G$ 相联系的自守层，实际特征值是系数域 $e$ 上的 Langlands 对偶群 $\checkG$ 的局部系统。

在自守一侧，$\GL_1=\Gm$ 的阿贝尔情形必须替换为其 $G$-类比，即
\[
\begin{array}{c|c|c|c}
\GL_1&K^\times\backslash\A_K^\times/\cO_K^\times&\Pic(X)&\Pic_{X/k}\\\hline
G&G(K)\backslash G(\A_K)/G(\cO_K)&?&?
\end{array}
\]
\begin{theorem}{Weil 一致化}{1.2.1}
在关于 $G$ 与 $X$ 的标准假设下，双陪集集合 $G(K)\backslash G(\A_K)/G(\cO_K)$ 自然等同于 $X$ 上主 $G$ 丛的同构类集合，也就是 $G$ 挠子的同构类集合。这些假设将在后面精确说明；这里所考虑的是曲线上的连通约化群。
\end{theorem}

与 $\Pic_{X/k}$ 不同，$G$ 丛的模对象通常不是概形。稳定子群随丛而变化，因此正确的对象是代数叠：它记住各点的自同构群。这些附加结构在单纯的双陪集集合中不可见，却对一般的 $G$ 至关重要。

\begin{proposition}{$\Bun_G$ 的定义}{1.2.2}
存在 $k$ 上的光滑 Artin 叠 $\Bun_G$，分类 $X$ 上的 $G$ 挠子族。在 $k$-点的同构类层面，
\[\Bun_G(k)\simeq G(K)\backslash G(\A_K)/G(\cO_K).\]
更精确的记号是 $\Bun_{G,X/k}$。
\end{proposition}
\srcpage{8}

Hecke 特征性质也需要推广。对 $G=\GL_1$，映射
\[\operatorname{add}:X\times_k\Pic_{X/k}\longrightarrow\Pic_{X/k},\qquad(x,L)\longmapsto L(x)\]
以典范方式在 $x$ 处修改线丛。对一般的 $G$ 挠子，在 $x$ 处并不存在唯一的典范修改，而是有许多种。因此，将映射 $\operatorname{add}$ 替换为整体 Hecke 对应
\origdiagram{p008-diagram.pdf}{0.5282}
它记录一个修改点，以及两个在该点之外彼此等同的 $G$ 丛。

假定已有 Hecke 特征层的一般定义，非分歧整体几何 Langlands 对应可以表述如下。
\begin{conjecture}{非分歧整体几何 Langlands}{1.2.3}
对 $X$ 上每个几何不可约的 $\checkG$-局部系统 $\sigma$，在 $\Bun_G$ 上都存在本质唯一的 Hecke 特征层 $\Aut_\sigma$，其特征值为 $\sigma$。
\end{conjecture}

这里，几何不可约是指：基变换到 $\bar k$ 后，相应 $\checkG$-局部系统不存在到 $\checkG$ 的真抛物子群的平坦约化；对于 $G=\GL_n$，这就是秩 $n$ 局部系统通常意义下的不可约性。如果 $\sigma$ 可约，预期得到的应是 $\Shv_c(\Bun_G)$ 中的 Hecke 特征对象，而不是心中的特征层；本质唯一性也不再成立。

对 $G=\GL_n$，存在性定理已经成立：Deligne 证明了 $n=1$ 的情形，Drinfeld 证明了 $n=2$ 的情形，Frenkel–Gaitsgory–Vilonen 证明了一般情形\cite{FGKV98}。在 $\ell$-进情形，对 $\Aut_\sigma$ 取 Frobenius 迹得到自守函数。因此，这个猜想给出函数域上从 Galois 数据到自守对象的几何构造；由 $\Aut_\sigma$ 得到的函数是尖点函数。

\section{非分歧整体几何 Langlands 等价}\label{sec:1.3}
\index{Langlands 等价!整体几何|hyperpage}\index{尖点|hyperpage}\index{奇异支撑!幂零|hyperpage}
\conjref{1.2.3}将一个自守层赋给一个不可约谱点。范畴形式则要求从整个谱模对象出发构造函子。以 $\LocSys_{\checkG}^{\mathrm{irred}}$ 表示分类 $X$ 上几何不可约 $\checkG$-局部系统族的代数几何对象。预期存在函子
\[\Coh(\LocSys_{\checkG}^{\mathrm{irred}})^\heartsuit\longrightarrow\Shv_c(\Bun_G)^\heartsuit,\qquad\delta_\sigma\longmapsto\Aut_\sigma,\]
而从理念上说，预期它与 $\Shv_c(\Bun_G)^\heartsuit$ 的尖点部分等价。

在 de Rham 情形，$\LocSys_{\checkG}^{\mathrm{irred}}$ 是光滑 Artin 叠，心层面的预期成为以下陈述。
\begin{conjecture}{尖点 de Rham 形式}{1.3.1}
存在典范的 $t$-正合等价
\[\mathbb L_G:\DMod_{\mathrm{coh}}(\Bun_G)_{\mathrm{cusp}}\simeq\Coh(\LocSys_{\checkG}^{\mathrm{irred}}),\]
使 $\mathbb L_G^{-1}$ 将不可约 $\checkG$-局部系统 $\sigma$ 处的摩天层映为以 $\sigma$ 为特征值的 Hecke 特征层 $\Aut_\sigma$。
\end{conjecture}

若将可约局部系统也包括进来，就要放弃 $t$-正合性，并在范畴层面工作。
\srcpage{9}
\begin{theorem}{朴素概括：完整 de Rham 形式}{1.3.2}
作为第一步近似，可以写出从谱一侧到自守一侧的等价
\[\mathbb L_G:\DMod_{\mathrm{coh}}(\Bun_G)\simeq\Coh(\LocSys_{\checkG}),\]
其中 $\checkG$-局部系统处的摩天层映为 Hecke 特征对象。这个式子只是一种指引性的概括：精确的非尖点 de Rham 表述，要将凝聚谱一侧替换为具有幂零奇异支撑的重整化范畴，例如 $\IndCoh_{\mathrm{Nilp}}(\LocSys_{\checkG})$。
\end{theorem}

这比经典 Langlands 对应高一个范畴层次：经典对象是函数与表示，几何对象则是层与 DG 范畴。更精确的现代表述将简化的凝聚谱一侧替换为适当的重整化范畴或具有幂零奇异支撑的范畴，例如 $\IndCoh_{\mathrm{Nilp}}(\LocSys_{\checkG})$\cite{AG15,GR24a}；但基本概括仍然是：$\Bun_G$ 上的自守范畴由 $\checkG$-局部系统上的谱范畴控制。这组讲授不讨论分歧整体理论。

\begin{transnote}{背景直觉}
「幂零奇异支撑」是 [AG15] 对谱一侧所作修正的核心。对 $\cF\in\QCoh(\LocSys_{\checkG})$，其奇异支撑 $\mathrm{SS}(\cF)\subset T^*\LocSys_{\checkG}$ 是余切丛里的闭锥；要求 $\mathrm{SS}(\cF)$ 落在幂零锥方向（即 $\LocSys^{\mathrm{nilp}}_{\checkG}$ 所指的那个锥）的对象构成全子范畴 $\IndCoh_{\mathrm{Nilp}}(\LocSys_{\checkG})$。它比 $\QCoh$ 小、比 $\Coh$ 大，并且在 Hecke 函子下稳定，因此才是与 $\DMod_{\mathrm{coh}}(\Bun_G)$ 相配的一侧。只保留此条件后，「几乎」所指的技术性缺口被填上，\S\ref{sec:1.3} 末尾的朴素整体形式才算精确。
[来源：nLab "geometric Langlands correspondence"，https://ncatlab.org/nlab/show/geometric+Langlands+correspondence（Last revised 2026-04-11；2026-11-17 访问）记 Arinkin–Gaitsgory 以 nilpotent ind-objects 替换 $\QCoh$ 并升级到稳定 $\infty$-范畴；并见 [AG15]（Selecta Math. 21 (2015), 1–199）的奇异支撑理论。]
\end{transnote}


\section{局部几何 Langlands 等价}\label{sec:1.4}
\index{Langlands 等价!局部几何|hyperpage}\index{环路群|hyperpage}\index{Whittaker 模型|hyperpage}
现在专注于 de Rham 情形。经典局部 Langlands 理论分类局部域 $K_v$ 上 $G(K_v)$ 的表示，它们构成一个 $1$-范畴。局部几何 Langlands 理论是其范畴化：它试图分类 $G(K_v)$ 的范畴表示，这些表示构成一个 $2$-范畴。

$G(K_v)$ 的几何体现是环路群 $LG$。在 $\mathbb C$ 上，它的 $\mathbb C$-点正是 $G(K_v)$ 的元素。定义 $LG\text{-}\mathrm{Mod}$ 为被幺半 DG 范畴 $\DMod(LG)$ 作用的 DG 范畴所构成的 $(\infty,2)$-范畴，其中幺半结构由卷积给出。若 $Y$ 是带有 $LG$ 作用的合理代数几何对象，则 $\DMod(Y)$ 是 $LG\text{-}\mathrm{Mod}$ 中的对象。

基本例子为 $\Gr_G$、$\Fl_G$ 和 $\Bun_G^{\infty;v}$：分别是仿射 Grassmann 簇、仿射旗 ind-簇，以及具有无限级结构的 $X$ 上 $G$ 挠子模 ind-叠，后者配备在 $v$ 的形式邻域 $D_v$ 上的平凡化。就复点而言，最后一个对象有双陪集描述
\[\prod_{w\ne v}G(\cO_w)\backslash G(\A_K)/G(K).\]

范畴 $\DMod(LG)$ 是它自身上的双模。对右作用取不变量，得到仍带有左 $LG$ 作用的 DG 范畴。于是，$\DMod(\Gr_G)$ 和 $\DMod(\Fl_G)$ 分别通过对右作用取 $L^+G$ 不变量和 $I$ 不变量得到。这里 $L^+G$ 是弧群，是 $G(\cO_v)$ 的几何对应物；$I$ 是 Iwahori 子群，是 $G(\cO_v)\times_{G(\kappa_v)}B(\kappa_v)$ 的几何对应物。

另一个重要对象是 Whittaker 模型
\[\Whit(LG)\in LG\text{-}\mathrm{Mod},\]
它通过相对于一般位置特征标 $\chi$ 取右 $LN$ 不变量而得到，其中 $N$ 是 Borel 子群 $B$ 的幂幺根基。还有一个关键对象 $\mathrm{KM}_{\mathfrak g}$，即临界水平下仿射 Lie 代数 $\widehat{\mathfrak g}$ 的表示所构成的 DG 范畴。

在 Galois 一侧，考虑穿孔圆盘 $D^\circ:=\Spec K_v$ 上 $\checkG$-局部系统的模问题 $\LocSys_{\checkG}(D^\circ)$。定义
\[\operatorname{QCohCat}(\LocSys_{\checkG}(D^\circ))\]
\srcpage{10}
为幺半范畴 $\QCoh(\LocSys_{\checkG}(D^\circ))$ 的模范畴所构成的 $(\infty,2)$-范畴，其中幺半结构为张量积。若 $Z\to\LocSys_{\checkG}(D^\circ)$ 是合理的代数几何对象，则 $\QCoh(Z)$ 就是这样的模范畴。

基本的谱一侧例子是
\[\LocSys_{\checkG}(D),\qquad\LocSys_{\checkG}^{\mathrm{nilp}}(D^\circ),\qquad\LocSys_{\checkG}(X^\circ),\]
分别分类圆盘 $D:=\Spec\cO_v$ 上的 $\checkG$-局部系统、穿孔圆盘上具有幂零奇点的 $\checkG$-局部系统，以及穿孔曲线 $X^\circ:=X\setminus\{v\}$ 上的 $\checkG$-局部系统。对约化群 $H$，有 $\LocSys_H(D)\simeq\mathrm{pt}/H$ 和 $\LocSys_H^{\mathrm{nilp}}(D^\circ)\simeq\cN/H$，其中 $\cN\subset\mathfrak h$ 是幂零锥。此外还考虑 $\LocSys_{\checkG}(D^\circ)$ 本身以及 oper 空间
\[\Op_{\checkG}(D^\circ).\]

局部猜想将上述自守一侧和谱一侧的例子配对；这就是局部几何 Langlands 的标准纲要形式\cite{FG05,FG07}。

\begin{conjecture}{局部几何 Langlands：纲要形式}{1.4.1}
存在几乎等价
\[LG\text{-}\mathrm{Mod}\simeq\operatorname{QCohCat}(\LocSys_{\checkG}(D^\circ)),\]
其中“几乎”掩盖了 $\QCoh$ 与凝聚或重整化变体之间的技术性不匹配：精确表述必须在两侧选取正确的凝聚部分。该等价具有以下对应：
\[
\begin{aligned}
\DMod(\Gr_G)&\longmapsto\QCoh(\LocSys_{\checkG}(D)),\\
\DMod(\Fl_G)&\longmapsto\QCoh(\LocSys_{\checkG}^{\mathrm{nilp}}(D^\circ)),\\
\DMod(\Bun_G^{\infty;v})&\longmapsto\QCoh(\LocSys_{\checkG}(X^\circ)),\\
\Whit(LG)&\longmapsto\QCoh(\LocSys_{\checkG}(D^\circ)),\\
\mathrm{KM}_{\mathfrak g}&\longmapsto\QCoh(\Op_{\checkG}(D^\circ)).
\end{aligned}
\]
\end{conjecture}

在列出的自守对象之间取 $\Hom$ 范畴，再与它们在谱一侧的像之间的 $\Hom$ 范畴比较，可以恢复许多熟悉的等价。在这些比较中，须将 $\DMod$ 替换为凝聚 $D$ 模，将 $\QCoh$ 替换为相应的凝聚谱范畴。
\begin{enumerate}
\item 根据 Bezrukavnikov–Finkelberg 的结果，导出几何 Satake 等价为
\[\End(\DMod(\Gr_G))\simeq\End(\QCoh(\LocSys_{\checkG}(D))).\]
\item Arkhipov–Bezrukavnikov–Ginzburg 定理为
\[
\begin{aligned}
&\Hom(\DMod(\Gr_G),\DMod(\Fl_G))\\
&\qquad\simeq\Hom(\QCoh(\LocSys_{\checkG}(D)),\QCoh(\LocSys_{\checkG}^{\mathrm{nilp}}(D^\circ))).
\end{aligned}
\]
\item 根据 Bezrukavnikov 关于仿射 Hecke 代数的两种几何实现之间等价的结果，
\[\End(\DMod(\Fl_G))\simeq\End(\QCoh(\LocSys_{\checkG}^{\mathrm{nilp}}(D^\circ))).\]
\item 有
\[
\begin{aligned}
&\Hom(\DMod(\Gr_G),\DMod(\Bun_G^{\infty;v}))\\
&\qquad\simeq\Hom(\QCoh(\LocSys_{\checkG}(D)),\QCoh(\LocSys_{\checkG}(X^\circ))).
\end{aligned}
\]
\srcpage{11}
这恢复了\thref{1.3.2}中的朴素整体形式；精确版本使用具有幂零奇异支撑的谱范畴。若将 $\DMod(\Gr_G)$ 和 $\QCoh(\LocSys_{\checkG}(D))$ 分别替换为 $\DMod(\Fl_G)$ 和 $\QCoh(\LocSys_{\checkG}^{\mathrm{nilp}}(D^\circ))$，就得到相应的驯顺分歧形式。
\item 根据 Frenkel–Gaitsgory–Vilonen 的结果\cite{FGKV98}，几何 Casselman–Shalika 公式为
\[
\begin{aligned}
&\Hom(\DMod(\Gr_G),\Whit(LG))\\
&\qquad\simeq\Hom(\QCoh(\LocSys_{\checkG}(D)),\QCoh(\LocSys_{\checkG}(D^\circ))).
\end{aligned}
\]
\item 根据 Arkhipov–Bezrukavnikov 的一个定理，
\[
\begin{aligned}
&\Hom(\DMod(\Fl_G),\Whit(LG))\\
&\qquad\simeq\Hom(\QCoh(\LocSys_{\checkG}^{\mathrm{nilp}}(D^\circ)),\QCoh(\LocSys_{\checkG}(D^\circ))).
\end{aligned}
\]
\item 以下等价称为 Feigin–Frenkel 定理\cite{Fre02}：
\[\Hom(\DMod(\Gr_G),\mathrm{KM}_{\mathfrak g})\simeq\Hom(\QCoh(\LocSys_{\checkG}(D)),\QCoh(\Op_{\checkG}(D^\circ))).\]
\item 根据 Raskin 的一个定理，
\[\Hom(\Whit(LG),\mathrm{KM}_{\mathfrak g})\simeq\Hom(\QCoh(\LocSys_{\checkG}(D^\circ)),\QCoh(\Op_{\checkG}(D^\circ))).\]
\end{enumerate}

\section{量子局部几何 Langlands 等价}\label{sec:1.5}
\index{Langlands 等价!量子局部几何|hyperpage}
扭曲 $D$ 模给出\conjref{1.4.1}中局部猜想的变形。预期的量子局部等价关于 $G$ 及其 Langlands 对偶群是对称的。
\begin{conjecture}{量子局部几何 Langlands}{1.5.1}
存在几乎等价
\[LG\text{-}\mathrm{Mod}_\kappa\simeq L\checkG\text{-}\mathrm{Mod}_{\check\kappa},\]
具有以下对应：
\[
\begin{aligned}
\DMod_\kappa(\Gr_G)&\longmapsto\DMod_{\check\kappa}(\Gr_{\checkG}),\\
\DMod_\kappa(\Fl_G)&\longmapsto\DMod_{\check\kappa}(\Fl_{\checkG}),\\
\DMod_\kappa(\Bun_G^{\infty;v})&\longmapsto\DMod_{\check\kappa}(\Bun_{\checkG}^{\infty;v}),\\
\Whit_\kappa(LG)&\longmapsto\mathrm{KM}_{\check{\mathfrak g},\check\kappa},\\
\mathrm{KM}_{\mathfrak g,\kappa}&\longmapsto\Whit_{\check\kappa}(L\checkG).
\end{aligned}
\]
\end{conjecture}

这个表格的对称性，就是 Langlands 对偶的量子版本。
